\cotefile{93}
\scanpage{3}
\margin{G? = correction de Grothendieck}
\margin{Pour une version un brin plus complète :\par
Lecture Notes 476 (Anvers \uline{IV}) p 53--73\par
see also: Tate, ibidem p 33--52 et spécialement p 33--39}

\begin{center}
COURBES ELLIPTIQUES : FORMULAIRE

 d'après J. Tate

mis au goût du jour par P. Deligne
\end{center}

1.\quad \uline{Formules générales}.

Dans ce formulaire, on appellera courbe de genre 1 sur un schéma de base $S$ un morphisme propre et plat de présentation finie
\[
p:E\longrightarrow S,
\]
muni d'une section $e$ contenue dans l'ouvert de lissité, dont les fibres géométriques sont des courbes réduites irréductibles de genre \uline{ar}ithmétique 1. Les fibres géométriques de $p$ sont donc soit :
\margin{G?}

 a) une courbe elliptique, i.e. une courbe propre, lisse et connexe de genre 1,

 b) une droite projective dont on a identifié deux points distincts (cubique dans $P_2$ à point double ordinaire),

 c) une droite projective dont on a identifié deux points infiniment voisins (cubique dans $P_2$ à singularité cuspidale).

On dira dans le cas a) (resp. b), c)) que la fibre est de type elliptique (resp. multiplicatif, resp. additif).

Si $E$ est une courbe de genre 1 sur $S$, on définit un faisceau inversible $\omega$ sur $S$ par la formule
\[
\omega=e^*\Omega^1_{E/S}.
\]
La section $e$ est un diviseur de Cortier relatif de $E$ sur $S$ ;
\margin{2}
